\section{The $\sl_2$-example in the chiral setting, at any level} \label{section:sl2-example}

\subsection{Definition of the localisation}

The universal affine vertex algebra $\V^k(\sl_2)$ is strongly generated by three vectors $e$, $h$ and $f$ whose only non-trivial OPE are given by
\[ \begin{multlined}
	h(z_1) h(z_2) \sim \frac{1}{(z_1 - z_2)^2} 2 k \id, \quad h(z_1) e(z_2) \sim \frac{1}{z_1 - z_2} 2 e(z_2) , \\ e(z_1) f(z_2) \sim \frac{1}{z_1 - z_2} h(z_2) + \frac{1}{(z_1 - z_2)^2} k \id \\ \text{and} \quad h(z_1) f(z_2) \sim - \frac{1}{z_1 - z_2} 2 f(z_2).
\end{multlined} \]
It is equipped with the grading \eqref{equation:grading-universal-affine} such that $e$ is in degree $0$, $h$ is in degree $1$ and $f$ in degree $2$.

By the general theory of \cite{borisov2001vertex,arakawa2026faithful}, the localised graded vertex algebra $\V^k(\sl_2)[e^{-1}]$ contains an additional generator $e^{-1}$, of degree $0$. Its associated field
\[ e^{-1}(z) \coloneqq \sum_{n \in \Z} e^{-1}_{(n)} \, z^{- n -1} \]
satisfies the following properties:
\begin{enumerate}
	\item All the coefficients $e^{-1}_{(n)}$ and $e_{(n)}$, for all $n$ in $\Z$, commute with each other, that is to say
	\[ e^{-1}(z_1) e(z_2) \sim 0.  \]
	
	\item The product
	\[ e(z) e^{-1}(z) \coloneqq \sum_{n \in \Z} \Big(\sum_{i + j = n} e_{(i)} \circ e^{-1}_{(j)} \Big) z^{- n -2} \]
	is well defined, meaning that the sum $\sum_{i + j = n} e_{(i)} \circ e^{-1}_{(j)} v$ is finite for any $v$ in $\V^k(\sl_2)[e^{-1}]$. This follows because of the vertex algebra grading: $e_{(i)} v$ or $e^{-1}_{(j)} v$ is zero for big values of $i$ and $j$.
	
	\item The identity 
	\[ e^{-1}(z) e(z) = e(z) e^{-1}(z) = \id \]
	is satisfied.
\end{enumerate}

As explained in \cite[Lemma 3.1 and before Lemma 3.3]{arakawa2026faithful}, the product $\alpha \coloneq e_{(-1)}$ is an invertible endomorphism of $\V^k(\sl_2)[e^{-1}]$, but in general
\[ \alpha^{-1} = (e_{(-1)})^{-1} \neq e^{-1}_{(-1)} \coloneqq (e^{-1})_{(-1)}, \] 
see \Cref{lemma:inverse-no}.

Everything in $\V^k(\sl_2)[e^{-1}]$ can be computed by considering sums of elements of the form
\[ \alpha^{-p}(v) \quad \text{for} \quad p \in \Z_{\geqslant 0} \quad \text{and} \quad v \in \V^k(\sl_2). \]
To compute the $e^{-1}_{(n)} v$, for $n$ in $\Z$ and $v$ in $\V^k(\sl_2)$, we must use the identity
\[ e(z) e^{-1}(z) v = v. \]
If $v$ is homogeneous of degree $\Delta$, then many terms equal zero and one gets
\[ \sum_{n \leqslant \Delta - 2} \Big(\sum_{\substack{i, j \leqslant \Delta - 1 \\ i + j = n}} e_{(i)} \circ e^{-1}_{(j)} v \Big) z^{- n -2} = v. \]

If we apply this identity to the vacuum $1$, we get:
\[ \sum_{n \leqslant - 2} \Big(\sum_{\substack{i, j \leqslant - 1 \\ i + j = n}} e_{(i)} \circ e^{-1}_{(j)} 1 \Big) z^{- n -2} = 1. \]
For $n = 2$, the identification gives
\[ e_{(-1)} \circ e^{-1}_{(j)} 1 = \NO{ e e^{-1}} = 1. \]
Therefore, the equality $e^{-1}_{(-1)} = ( e_{(-1)})^{-1}$ holds on the commutative vertex subalgebra $\big(\V^k(\sl_2)[e^{-1}]\big)_0$ which is generated by $e$ and $e^{-1}$ (it is the lowest part of the increasing Li filtration): $e^{-1}$ is the inverse of $e$ in the usual associative sense.

The translation operator $\partial$ satisfies a Leibniz rule for the normally ordered product:
\[ \partial(\NO{ab}) = \NO{\partial(a) b} + \NO{a \partial(b)}. \]
Therefore, we get
\[ \partial(e^{-1}) = - \NO{e \NO{e \partial(e)}} = - e^{-2} \partial e, \]
where we can drop the normally ordered product because everything commutes.

\emph{In general, if an expression only involves $e$, $e^{-1}$ and their derivatives, we omit the $\NO{}$ to denote the product because everything will commute.}

\subsection{Operator product expansions} \label{subsection:ope-localised-sl2}

We compute the OPE between $e^{-1}$ and the other generators manually, by using the machinery described above.

\begin{proposition} \label{proposition:relations-localisation-sl2}
	The localised universal affine vertex algebra $\V^k(\sl_2)[e^{-1}]$ is strongly generated by four vectors $e$, $e^{-1}$, $h$ and $f$ whose only non-trivial OPE are given by
	\[ \begin{multlined}
		h(z_1) h(z_2) \sim \frac{1}{(z_1 - z_2)^2} 2 k \id, \quad h(z_1) e(z_2) \sim \frac{1}{z_1 - z_2} 2 e(z_2) , \\ e(z_1) f(z_2) \sim \frac{1}{z_1 - z_2} h(z_2) + \frac{1}{(z_1 - z_2)^2} k \id, \\
		h(z_1) f(z_2) \sim - \frac{1}{z_1 - z_2} 2 f(z_2),\quad
		e^{-1}(z_1) h(z_2) \sim  \frac{1}{z_1 - z_2} 2 e^{-1}(z_2), \\
		\text{and} \quad e^{-1}(z_1) f(z_2) \sim  \frac{1}{z_1 - z_2} \Big( 2 (k - 1) \big((e^{-1})^3 \partial e\big)(z_2) - \NO{e^{-1} e^{-1} h}(z_2)\Big) \\
		+ \frac{1}{(z_1 - z_2)^2} \Big( - (k + 2) (e^{-1})^2(z_2)\Big).
	\end{multlined} \]
	and satisfying the additional relation
	\[ e^{-1}(z) e(z) = e(z) e^{-1}(z) = \id. \]
\end{proposition}

\begin{proof}[Computation of the OPE between $e^{-1}$ and $h$]
	The operator product expansion of $e$ and $h$ is
	\[ e(z_1) h(z_2) \sim \frac{1}{z_1 - z_2} (- 2 e(z_2)), \]
	therefore $e_{(0)} h = - 2 e$ and all other products are zero. The element $h$ is of degree $1$, so the product $e^{-1}_{(n)} h$ is of degree $0 + 1 - n - 1 = -n$, that is a negative number for $n > 0$. So only the product $e^{-1}_{(0)} h$ may be non-zero. Consider the identity
	\[ \sum_{n \leqslant- 1} \Big(\sum_{\substack{i, j \leqslant 0 \\ i + j = n}} e_{(i)} \circ e^{-1}_{(j)} h \Big) z^{- n -2} = h. \]
	For $n = -1$, we get the equality
	\[ e_{(-1)} \circ e^{-1}_{(0)} h + e_{(0)} \circ e^{-1}_{(-1)} h = e_{(-1)} \circ e^{-1}_{(0)} h +  e^{-1}_{(-1)} \circ e_{(0)} h  = 0. \]
	By using the equality
	\[ e^{-1}_{(-1)} h =  e^{-1}_{(-1)} \circ e_{(0)} h = - 2 e^{-1}_{(-1)} e = - 2 , \]
	we get
	\[ e^{-1}_{(0)} h = 2 e_{(-1)} 1 = 2 e^{-1}. \]
\end{proof}

The OPE between $e^{-1}$ and $f$ is slightly more involved. We need the following lemma.

\begin{lemma} \label{lemma:inverse-no}
	The following equality holds:
	\[ (e_{(-1)})^{-1} (\NO{e^{-1} h}) - \NO{e^{-1} e^{-1} h} = 4 (e^{-1})^3 \partial e. \]
	This proves in particular that $(e_{(-1)})^{-1} \neq e^{-1}_{(-1)} = \NO{e^{-1} \bullet}$.
\end{lemma}

\begin{proof}
	Localisation imposes the equality:
	\[ e_{(0)} e^{-1}_{(-2)} (\NO{e h^{-1}}) + e_{(-1)} e^{-1}_{(-1)} (\NO{e h^{-1}}) + e_{(-2)} e^{-1}_{(0)} (\NO{e h^{-1}}) = \NO{e h^{-1}}. \]
	By the Borcherds identity \eqref{equation:borcherds}:
	\[ \begin{aligned}
		e_{(0)}(\NO{e^{-1} h^{-1}}) & = \NO{e^{-1} (e_{(0)} h)} = - 2 \NO{e^{-1}e } = - 2, \\
		e^{-1}_{(0)}(\NO{e^{-1} h^{-1}}) & = \NO{e^{-1} (e^{-1}_{(0)} h)} = 2 (e^{-1})^2.
	\end{aligned} \]
	Moreover,
	\[ \begin{aligned}
		e_{(0)} e^{-1}_{(-2)} (\NO{e h^{-1}}) & = e^{-1}_{(-2)} e_{(0)} (\NO{e h^{-1}}) = - 2 e^{-1}_{(-2)}1 = - 2 \partial(e^{-1}) = 2 (e^{-1})^2 \partial e, \\
		e_{(-2)} e^{-1}_{(0)} (\NO{e h^{-1}}) & =e_{(-2)} e^{-1}_{(0)}   (\NO{e h^{-1}}) = 2 (e^{-1})^2 \partial e.
	\end{aligned} \]
	From this we deduce that
	\[ 4 e^{-1}e^{-1} \partial e + e_{(-1)} e^{-1}_{(-1)} (\NO{e h^{-1}})  = (\NO{e h^{-1}}). \]
	By applying $(e_{(-1)})^{-1}$, we get the identity stated in the lemma.
\end{proof}

Now we can finish the proof of \Cref{proposition:relations-localisation-sl2}.

\begin{proof}[Computation of the OPE between $e^{-1}$ and $f$]
	Because $e^{-1}$ is of degree $0$ and $f$ is of degree $2$, only the products $e^{-1}_{(0)}f$ and $e^{-1}_{(1)}f$ may not be zero. We have $e_{(0)} f = h$ and $e_{(1)} f = k 1$, all other products are zero.
	
	Localisation imposes the equality:
	\[ e_{(1)} e^{-1}_{(-1)} f + e_{(0)} e^{-1}_{(0)} f + e_{(-1)} e^{-1}_{(1)} f = 0. \]
	Direct computations give:
	\[ \begin{aligned}
		e_{(1)} e^{-1}_{(-1)} f & = e^{-1}_{(-1)} e_{(1)} f = k e^{-1}, \\
		e_{(0)} e^{-1}_{(0)} f & = e^{-1}_{(0)} e_{(0)} f = e^{-1}_{(0)} h = 2 e^{-1}.
	\end{aligned}  \]
	So, we get the equality
	\[ k e^{-1} + 2 e^{-1} + e_{(-1)} e^{-1}_{(1)} f = 0, \]
	that is to say
	\[ e^{-1}_{(1)} f  = - (k + 2) (e^{-1})^2, \]
	since $(e_{(-1)})^{-1} e^{-1} = (e^{-1})^2$.
	
	Localisation also imposes the identity
	\[ e_{(1)} e^{-1}_{(-2)} f + e_{(0)} e^{-1}_{(-1)} f + e_{(-1)} e^{-1}_{(0)} f + e_{(-2)} e^{-1}_{(1)} f = 0. \]
	By similar computations as before, we get:
	\[ e^{-1}_{(0)} f = (2 k + 2)(e^{-1})^3 \partial e - (e_{(-1)})^{-1} (\NO{e^{-1} h}). \]
	But \Cref{lemma:inverse-no} gives
	\[ \begin{aligned}
		e^{-1}_{(0)} f & = (2 k + 2) (e^{-1})^3 \partial e - (\NO{e^{-1} e^{-1} h} + 4 (e^{-1})^3 \partial e) \\
		& = 2 (k - 1) (e^{-1})^3 \partial e - \NO{e^{-1} e^{-1} h}.
	\end{aligned} \]
\end{proof}



The relation provided by \Cref{proposition:relations-localisation-sl2} can be implemented in Mathematica by using the package \emph{OPEdefs} \cite{thielemans1991mathematica}.



\subsection{The affine W-algebra is a vertex subalgebra of the localisation}

We recall the following well-known fact, see for example \cite[Section 15.3.2]{frenkel2004vertex}.

\begin{theorem} \label{theorem:w-algebra-sl2}
	If $k = -2$, the W-algebra $W^{-2}(\sl_2, f)$ is isomorphic to the centre $\operatorname{Z}(\V^k(\sl_2))$ of the universal affine vertex algebra. This centre is freely generated by
	\[ \widehat{c} \coloneqq \frac{1}{2} \NO{h h} + 2 \NO{e f} - \partial h. \]
	
	If $k \neq -2$, then the W-algebra is isomorphic to the Virasoro vertex algebra of central charge 
	\[ c(k) \coloneqq 1 - \frac{6 (k + 1)^2}{k + 2}. \]
	That means that $\W^k(\sl_2, f)$ is freely generated by one element $\omega$ satisfying the OPE
	\begin{equation} \label{equation:virasoro} \omega(z_1) \omega(z_2) \sim \frac{1}{z_1 - z_2} \partial \omega(z_2) + \frac{2}{(z_1 - z_2)^2} \omega(z_2) + \frac{1}{(z_1 - z_2)^4} \frac{c(k)}{2} \id.  \end{equation}
\end{theorem}

\begin{proposition} \label{proposition:invariant-generator-sl2}
	The invariant subalgebra $\big(\V^k(\sl_2)[e^{-1}]\big)^{\J_\infty(\C h \oplus \C e)}$ is freely generated by the element
	\[ \begin{multlined}
		\widehat{c}_k \coloneqq \frac{1}{2} \NO{h h} + 2 \NO{e f} + (1 + k) \partial h \\ 
		+ (k + 2) \Big(\frac{3 k - 4}{2} (e^{-1})^2 (\partial e)^2 -  \NO{e^{-1} (\partial e) h} - (1+k) \NO{e^{-1} \partial^2 e}\Big).
	\end{multlined}  \]
	\begin{itemize}
		\item If $k = -2$, then $\widehat{c}_{-2} = \widehat{c}$ is the strong generator of the centre $\operatorname{Z}(\V^k(\sl_2))$.
		
		\item If $k \neq -2$, then $\omega \coloneqq \frac{1}{2(k + 2)} \widehat{c}_k$ satisfies the OPE \eqref{equation:virasoro}.
	\end{itemize}
\end{proposition}

\begin{proof}
	The vertex algebra grading and the eigenvalues of the endomorphism $h_{(0)}$ define a bi-grading:
	\[ \begin{multlined}
		\V^k(\sl_2)[e^{-1}] = \bigoplus_{\substack{\Delta \in \Z_{\geqslant 0} \\ \delta \in \Z}} \V^k(\sl_2)[e^{-1}](\Delta, \delta), \\
		\text{where} \quad \V^k(\sl_2)[e^{-1}](\Delta, \delta) \coloneqq \{ a \in \V^k(\sl_2)[e^{-1}](\Delta) \quad | \quad h_{(0)} a = \delta a \}.
	\end{multlined}  \]
	Each piece is finite-dimensional. 
	
	\Cref{theorem:affine-w-algebra-coset} and \Cref{theorem:w-algebra-sl2} state that $ \big(\V^k(\sl_2)[e^{-1}]\big)^{\J_\infty(\C h \oplus \C e)}$ is freely generated by an element $\widehat{c}_k $ that must be in degree $2$ for the vertex algebra grading. This element is in degree $0$ for the action of $h$ because it is $h_{(0)}$-invariant by definition. As a vector space,
	\[ \V^k(\sl_2)[e^{-1}](2, 0) = \operatorname{Span}\big(\NO{h h},  \NO{e f}, \NO{\partial h}, (e^{-1})^2 (\partial e)^2, \NO{e^{-1} (\partial e) h}, \NO{e^{-1} \partial^2 e}\big)  \]
	is $6$-dimensional.
	
	The subspace 
	\[ \Big(\V^k(\sl_2)[e^{-1}](2, 0)\Big)^{\J_\infty(\C h \oplus \C e)}  \] 
	can be computed by brute force computation. Taking a generic element in $\V^k(\sl_2)[e^{-1}]$ and imposing the commutation with $h$ and $e$ gives a set of linear equations (by using \emph{OPEdefs}). Solving this equation leads to a one dimensional vector space
	\[ \Big(\V^k(\sl_2)[e^{-1}](2, 0)\Big)^{\J_\infty(\C h \oplus \C e)}  = \C \widehat{c}_k, \]
	where $\widehat{c}_k$ is given by the formula in \Cref{proposition:invariant-generator-sl2}.
	
	If $k = -2$, the statement is clear. For $k \neq -2$, we can write
	\[ \omega_1 = \omega + \omega_2, \]
	where
	\[ \begin{aligned}
		\omega_1 & \coloneqq \frac{1}{2(k + 2)} \Big(\frac{1}{2} \NO{h h} + 2 \NO{e f} - \partial h\Big) + \frac{1}{2} \partial h, \\
		\omega_2 & \coloneqq - \Big(\frac{3 k - 4}{2} (e^{-1})^2 (\partial e)^2 -  \NO{e^{-1} (\partial e) h} - (1+k) \NO{e^{-1} \partial^2 e}\Big).
	\end{aligned} \]
	By construction of $\omega$ as an element commuting with $e$ and $h$, we get that
	\[ \omega(z_1) \omega_2(z_2) \sim 0. \]
	
	According to \cite[Example 1.4]{kac2003quantum}, $\omega_1$ is the $h$-twisted Sugawara conformal vector in $\V^k(\sl_2)$, meaning that we have the OPE:
	\[ \omega_1(z_1) \omega_1(z_2) \sim \frac{1}{z_1 - z_2} \partial \omega_1(z_2) + \frac{2}{(z_1 - z_2)^2} \omega_1(z_2) + \frac{1}{(z_1 - z_2)^4} \frac{1}{2} \frac{(-3 k)(2 k + 3)}{k + 2} \id. \] 
	A brute force computation with \emph{OPEdefs} gives 
	\[ \omega_2(z_1) \omega_2(z_2) \sim \frac{1}{z_1 - z_2} \partial \omega_2(z_2) + \frac{2}{(z_1 - z_2)^2} \omega_2(z_2) + \frac{1}{(z_1 - z_2)^4} \frac{1}{2} 2 \id. \]
	
	For $n \geqslant 0$, we have $\omega_{(n)} \omega_2 = \omega_2{}_{(n)} \omega = 0$. Therefore
	\[ \omega_1{}_{(n)} \omega_1 = \omega_{(n)} \omega + \omega_2{}_{(n)} \omega_2. \]
	Therefore we get the OPE
	\[ \omega(z_1) \omega(z_2) \sim \frac{1}{z_1 - z_2} \partial \omega(z_2) + \frac{2}{(z_1 - z_2)^2} \omega_1(z_2) + \frac{1}{(z_1 - z_2)^4} \frac{1}{2} c(k) \id, \]
	where
	\[ c(k) = \frac{(-3 k)(2 k + 3)}{k + 2} - 2 = 1 - \frac{6 (k + 1)^2}{k + 2}. \]
\end{proof}


\Cref{theorem:decomposition-chiral} gives an isomorphism
\[ \widehat{\delta} \colon  \V^k(\C h \oplus \C e) \otimes  \big(\V^k(\sl_2)[e^{-1}]\big)^{\J_\infty(\C h \oplus \C e)} \overset{\sim}{\longrightarrow} \big(\V^k(\sl_2)[e^{-1}]\big)^{\J_\infty(\C h \oplus \C e)}. \]
In this case, we can retrieve the inverse reduction \eqref{equation:inverse-reduction} by computing $\widehat{\delta}^{-1}(f)$.

\begin{proposition} \label{proposition:computation-f}
	The following equality
	\[ \begin{multlined}
		f = \NO{e^{-1} \frac{1}{2} \widehat{c}_k} + \frac{(k - 2)(3 k -4)}{4} (e^{-1})^3 (\partial e)^2 \\ + \frac{k - 2}{2} \NO{e^{-1} e^{-1} (\partial e) h} + \frac{(k^2 + k - 4)}{2} \NO{e^{-1} h h} - \frac{k + 1}{2} \NO{e^{-1} \partial h}.
	\end{multlined} \]
	gives an expression of $\widehat{\delta}^{-1}(f)$ in term of $e$, $e^{-1}$, $h$ and $\widehat{c}_k$. 
\end{proposition}

\begin{proof}
	The Borcherds identity \eqref{equation:borcherds} gives:
	\[ f = \NO{e^{-1} e f} + (3k - 4) (e^{-1})^3 (\partial e)^2 - 2 \NO{e^{-1} e^{-1} (\partial e) h} - (k + 3) (e^{-1})^2 \partial^2 e. \]
	Therefore:
	\[ \begin{multlined}
		f = \NO{e^{-1} \frac{1}{2} \widehat{c}_k}  - \NO{e^{-1} \frac{1}{2} \bigg(\frac{1}{2} \NO{h h} + (1 + k) \partial h + \\ (k + 2) \Big(\frac{3 k - 4}{2} (e^{-1})^2 (\partial e)^2  -  \NO{e^{-1} (\partial e) h} - (1+k) \NO{e^{-1} \partial^2 e}\Big)\bigg)}. \\  
		+ (3k - 4) (e^{-1})^3 (\partial e)^2 - 2 \NO{e^{-1} e^{-1} (\partial e) h} - (k + 3) (e^{-1})^2 \partial^2 e. 
	\end{multlined} \]
	After simplification, we get the desired formula.
	
\end{proof}
